\documentclass{subfiles}
\begin{document}
    We now consider an interesting result about UD codes: the Kraft-McMillan inequality.
        Formally, the Kraft-McMillan theorem stastes that a code \(\C*\) over an alphabet 
        of size \(d\), with of codewords lengths \(l_{1}, \ldots, l_{n}\) is UD, if and only if 
        \[
            \sum_{i = 1}^{n}d^{-l_{i}} \le 1.
        \]

    The inequality provides therefore a necessary and sufficient condition for a
        code to be UD. 

    \begin{proof*}
        We prove the result by construction; that is, 
            we show that if a set of lengths satisfies the inequality,
            then we can build a tree whose paths from the root to the leaves 
            define codewords of a prefix-free code.

        \begin{description}
            \item [\(\impliedby\)] Easely seen with any UD code.
            \item [\(\implies\)] Let \(l_{1}, \ldots, l_{n}\) be such that 
                \[
                    \sum_{i = 1}^{n} d^{-l_{i}} \le 1.
                \]
                Denote by \(n_{j}\) the number of codewords of length \(j\).
                Let \(L\) be the maximum length.
                Then 
                \[
                    \sum_{j = 1}^{L} n_{j} = n
                \]
                since there are \(n\) codewords. Hence, we have 
                \[
                    \sum_{i = 1}^{n} d^{-l_{i}} = \sum_{j = 1}^{L} n_{j}d^{-j} \le 1.
                \]
                By doing so, we build a \(d\)-ary tree which has \(n_{1}\) 
                leaves at depth 1, \(n_{2}\) leaves at depth 2, etc.
                It's easely seen that this produces a prefix-free code.\qed
        \end{description}
    \end{proof*}

    \begin{remark*}
        If the set of lengths is such that the sum is strictly less then one,
            then some paths in the tree will correspond to no codeword.
    \end{remark*}
\end{document}
